%! Trigonometric Equations and Inequalities — Unit 3, Lesson 3.10 — Lesson Plan
\documentclass[10pt]{article}
\usepackage{precalculus-article}
\usepackage{precalculus-boxes}
\usepackage{graphicx}
\graphicspath{{images/}}
\newcommand{\TallMath}[1]{$\displaystyle #1\rule[-1.4em]{0pt}{3.2em}$}
\newcommand{\CourseName}{AP Precalculus}
\newcommand{\SchoolYear}{2026--2027}
\newcommand{\MeetingLength}{55 minutes}
\newcommand{\UnitNumberName}{Unit 3: Trigonometric and Polar Functions \quad}
\newcommand{\LessonNumberName}{Lesson 3.10: Trigonometric Equations and Inequalities}

\begin{document}

%----------------------------------------------------------------------
% Title Block
%----------------------------------------------------------------------
\begin{center}
  {\Large\bfseries \CourseName: \SchoolYear} \\[4pt]
  {\small\bfseries \UnitNumberName \LessonNumberName \quad (3--4 instructional periods, \MeetingLength\ each)}
\end{center}

%----------------------------------------------------------------------
% Primary Objective
%----------------------------------------------------------------------
\begin{tcolorbox}[colback=sky, colframe=navy, boxrule=0.9pt, arc=2mm,
    left=2mm, right=2mm, top=1.5mm, bottom=1.5mm]
  \textbf{Primary Objective:} Students will solve equations and inequalities
  involving trigonometric functions by isolating the trig function, applying
  inverse trigonometric functions to find the principal value, and using
  periodicity to generate the general solution. In contextual scenarios,
  students will identify the implied domain restriction and select only
  the solutions that are valid in context.
  \textit{(Big Idea: Functions)}
\end{tcolorbox}

\vspace{0.08in}

%----------------------------------------------------------------------
% Priority Ideas & Skills
%----------------------------------------------------------------------
\begin{skillbox}[Priority Ideas \& Skills]{goldbox}
\begin{tabularx}{\linewidth}{@{} p{0.46\linewidth} X @{}}
\textbf{AP Skills} & \textbf{Key Understandings (EKs)} \\
\midrule
\textbf{1.A} \textit{Solve equations/inequalities analytically} ---
  Isolate a trigonometric function, apply an inverse trig function to
  find the principal value, then use periodicity to write the general
  solution.
&
\textbf{EK 3.10.A.1}: Inverse trigonometric functions are useful in
solving equations and inequalities involving trigonometric functions,
but solutions may need to be modified due to domain restrictions of
the inverse function. \\[4pt]

\textbf{2.A} \textit{Identify information from representations} ---
  Read solution intervals from a contextual function model (graph or
  formula) and connect them to the question being asked.
&
\textbf{EK 3.10.A.2}: Because trigonometric functions are periodic,
there are often infinitely many solutions to trigonometric equations
when the domain is unrestricted. \\[4pt]

\textbf{3.B} \textit{Apply numerical results} ---
  Interpret solutions in context: select values of $t$ (or $\theta$)
  that lie within the implied domain of the scenario.
&
\textbf{EK 3.10.A.3}: In trigonometric equations and inequalities
arising from a contextual scenario, there is often an implied domain
restriction that limits the number of valid solutions. \\
\end{tabularx}
\end{skillbox}

\vspace{0.08in}

%----------------------------------------------------------------------
% Vocabulary
%----------------------------------------------------------------------
\begin{skillbox}[{Vocabulary, Concepts, \& Theorems}]{greenbox}
\begin{tabularx}{\linewidth}{@{} p{0.28\linewidth} X @{}}
\textbf{Term} & \textbf{Definition / Note} \\
\midrule
trigonometric equation &
  An equation in which the unknown appears inside a trigonometric
  function ($\sin$, $\cos$, $\tan$, etc.). \\[4pt]
principal solution &
  The solution obtained by applying an inverse trig function directly;
  restricted to the range of that inverse function. \\[4pt]
general solution &
  All solutions to a trigonometric equation, expressed using the period
  (e.g.\ $\theta = \tfrac{\pi}{6} + 2\pi k$, $k \in \mathbb{Z}$). \\[4pt]
domain restriction &
  A limit on the variable imposed by the context of the problem
  (e.g.\ $0 \le t \le 24$ for a 24-hour tide model). \\
\end{tabularx}
\end{skillbox}

\vspace{0.08in}

%----------------------------------------------------------------------
% Activate Prior Knowledge
%----------------------------------------------------------------------
\begin{skillbox}[Activate Prior Knowledge \& Spiral Review]{sky}
\begin{tabularx}{\linewidth}{@{}X@{}}
\textbf{Spiral topics activated during Warm-Up (first 8 minutes):}
\begin{itemize}[leftmargin=1.4em, itemsep=2pt]
  \item \textbf{Inverse trig values}: Evaluate $\arcsin\!\bigl(\tfrac{\sqrt{2}}{2}\bigr)$
    and $\arccos\!\bigl(-\tfrac{1}{2}\bigr)$ --- confirms students recall the
    ranges of $\arcsin$ and $\arccos$ before the lesson relies on them.
  \item \textbf{Unit circle / reference angles}: Given $\sin\theta = \tfrac{1}{2}$,
    find all $\theta \in [0, 2\pi)$ --- prerequisite for generating the second
    solution family beyond the principal value.
  \item \textbf{Isolating a trig function}: Solve $2\sin x + 1 = 0$ by isolating
    $\sin x$ (no inverse trig yet) --- reinforces the algebraic manipulation that
    opens every trig equation.
\end{itemize}
\end{tabularx}
\end{skillbox}

\vspace{0.08in}

%----------------------------------------------------------------------
% Hook
%----------------------------------------------------------------------
\begin{skillbox}[Hook --- Entry Event]{sky}
\textbf{Scenario:} Display the graph of $h(t) = 3\sin\!\bigl(\tfrac{\pi}{6}t\bigr) + 5$,
representing tidal height (in feet) over a 24-hour period.

\textbf{Discussion prompts:}
\begin{enumerate}[leftmargin=1.4em, itemsep=3pt]
  \item \textit{``A harbor requires at least 7 feet of water for a ship to
    dock. During what hours today can the ship safely enter?''}
  \item \textit{``If the tide is exactly 7 feet, how many times does that
    happen in one 24-hour period? Could it happen again tomorrow?''}
  \item \textit{``How is this question different from `find the maximum tide'?
    What new mathematical tool do we need?''}
\end{enumerate}

\textbf{Key tension to surface:} Reading a graph gives approximate answers.
Finding exact times requires solving $3\sin\!\bigl(\tfrac{\pi}{6}t\bigr) + 5 = 7$
analytically. Today's tools --- inverse trig functions plus periodicity ---
give exact solutions and also reveal the full interval where the inequality
$h(t) \ge 7$ holds.

\textbf{Transition:} ``By the end of this lesson you will solve that equation
exactly, write all solutions as a general formula, and determine precisely
which hours the ship can enter.''
\end{skillbox}

\vspace{0.08in}

%----------------------------------------------------------------------
% Lesson
%----------------------------------------------------------------------
\begin{skillbox}[Lesson --- Instruction Sequence]{sky}
\begin{multicols}{2}

\textbf{Step 1 --- Infinitely Many Solutions} (10 min)

Motivate EK 3.10.A.2 with $\sin\theta = \tfrac{1}{2}$.
On the unit circle, mark both solutions in $[0, 2\pi)$:
$\theta = \tfrac{\pi}{6}$ and $\theta = \tfrac{5\pi}{6}$.
Ask: ``Can we add $2\pi$ and get another solution?''
Yes --- leading to infinitely many in an unrestricted domain.
Write the \emph{general solution} in two families:
$\theta = \tfrac{\pi}{6} + 2\pi k$ and $\theta = \tfrac{5\pi}{6} + 2\pi k$,
$k \in \mathbb{Z}$.
Emphasize: the $2\pi k$ comes from the period of sine; each trig function
contributes its own periodicity to the general solution.

\columnbreak

\textbf{Step 2 --- Inverse Trig Procedure} (12 min)

Teach the four-step algorithm (EK 3.10.A.1) using $2\cos\theta - 1 = 0$:
\begin{enumerate}[leftmargin=1.2em, itemsep=1pt, topsep=2pt]
  \item Isolate: $\cos\theta = \tfrac{1}{2}$.
  \item Principal value: $\theta = \arccos\!\bigl(\tfrac{1}{2}\bigr) = \tfrac{\pi}{3}$.
  \item Second solution in $[0,2\pi)$: $\theta = 2\pi - \tfrac{\pi}{3} = \tfrac{5\pi}{3}$
    (cosine is positive in Q1 and Q4, so reflect over the $x$-axis).
  \item General solution: $\theta = \tfrac{\pi}{3} + 2\pi k$ or
    $\theta = \tfrac{5\pi}{3} + 2\pi k$, $k \in \mathbb{Z}$.
\end{enumerate}
Discuss: the second solution in Step 3 depends on which quadrants give the
correct sign of the trig function --- this is why knowing the unit circle is
essential.

\end{multicols}

\begin{multicols}{2}

\textbf{Step 3 --- Contextual Domain Restriction} (12 min)

Return to the hook model $h(t) = 3\sin\!\bigl(\tfrac{\pi}{6}t\bigr) + 5$.
Solve the equation $h(t) = 7$:
\begin{enumerate}[leftmargin=1.2em, itemsep=1pt, topsep=2pt]
  \item Isolate: $\sin\!\bigl(\tfrac{\pi}{6}t\bigr) = \tfrac{2}{3}$.
  \item Let $u = \tfrac{\pi}{6}t$; solve $\sin u = \tfrac{2}{3}$.
  \item Principal: $u = \arcsin\!\bigl(\tfrac{2}{3}\bigr)$; second: $u = \pi - \arcsin\!\bigl(\tfrac{2}{3}\bigr)$.
  \item Solve for $t$; apply domain $t \in [0, 24]$ to select valid values.
\end{enumerate}
Then solve $h(t) \ge 7$ by identifying the interval between the two
boundary solutions where sine exceeds $\tfrac{2}{3}$. Emphasize EK 3.10.A.3:
context ($0 \le t \le 24$) means we keep only the solutions in that window.

\columnbreak

\textbf{Step 4 --- Trig Inequalities} (10 min)

Generalize the structure for $\sin\theta \ge c$ ($|c| \le 1$):
\begin{itemize}[leftmargin=1.2em, itemsep=2pt, topsep=2pt]
  \item Find the two boundary solutions in one period.
  \item Sine $\ge c$ on the arc between those values (the ``upper'' arc
    of the unit circle).
  \item Extend using $2\pi k$ if the domain is unrestricted.
\end{itemize}
Briefly contrast $\cos\theta \le c$ (solution on the ``lower'' arc) and
$\tan\theta > c$ (solution spanning half a period).
Close: \textit{``Why does the contextual domain sometimes give you fewer
solutions than the general formula?''}

\end{multicols}

\smallskip
\textbf{Pacing note:} Steps 1--2 in Periods 1--2; Steps 3--4 + activity
in Periods 3--4.
\end{skillbox}

\vspace{0.08in}

%----------------------------------------------------------------------
% Active Monitoring
%----------------------------------------------------------------------
\begin{skillbox}[Active Monitoring]{redbox}
\begin{itemize}[leftmargin=1.4em, itemsep=3pt]
  \item \textbf{Missing second solution}: Students often find only the principal
    value and stop. Prompt: \textit{``Sine (or cosine) takes this value in two
    quadrants per period --- which other angle gives the same value?''}
  \item \textbf{Forgetting $+2\pi k$}: After finding two solutions in one
    period, students omit the general form. Ask: \textit{``How many solutions
    exist if the domain is all real numbers?''}
  \item \textbf{Domain restriction errors}: Students list general-solution values
    without filtering. Prompt: \textit{``Does this $t$ value fall within the
    24-hour context? Is it physically meaningful?''}
  \item \textbf{Inequality direction}: Reinforce that isolating the trig function
    (dividing by a positive coefficient) does not flip the inequality sign;
    direction flips only when dividing by a negative.
\end{itemize}
\end{skillbox}

\vspace{0.08in}

%----------------------------------------------------------------------
% Group Work & Differentiation
%----------------------------------------------------------------------
\begin{skillbox}[Group Work \& Differentiation]{redbox}
Students work in groups of 3--4 on the tiered activity (assign by teacher
or allow self-selection with guidance).

\begin{itemize}[leftmargin=1.4em, itemsep=4pt]
  \item \textbf{Tier R (Remediate)}: General solution formulas are provided.
    Students substitute specific integers $k = 0, 1, 2, -1, -2$ to list all
    solutions in $[0, 4\pi]$, then circle those in a specified sub-interval.
    Step-by-step scaffold included.
  \item \textbf{Tier A (Apply)}: Students solve three trig equations without
    scaffolding: one straightforward, one requiring a quadratic approach
    ($u = \sin\theta$), and one contextual problem requiring a domain
    restriction. Students find all solutions in the given interval.
  \item \textbf{Tier E (Extend)}: Students solve a trig inequality using a
    table of pre-computed values of $f(\theta) = 2\cos\theta - 1$ to identify
    where $f(\theta) > 0$, then write the solution as a union of intervals.
    One contextual problem requires selecting and interpreting valid solutions
    in writing.
\end{itemize}
\end{skillbox}

\vspace{0.08in}

%----------------------------------------------------------------------
% Individual Work & Assessment
%----------------------------------------------------------------------
\begin{skillbox}[Individual Work \& Assessment]{redbox}
Exit ticket administered independently (final 5 minutes of Period 4).

\begin{enumerate}[leftmargin=1.4em, itemsep=4pt]
  \item Solve $\sin\theta = -\dfrac{\sqrt{3}}{2}$ for all $\theta \in [0, 2\pi]$.
  \item Explain in one or two sentences why a trigonometric equation has
    infinitely many solutions when the domain is all real numbers.
\end{enumerate}

\textbf{Data use:} Sort exit tickets into three groups --- correct general
solution (both families), correct principal value only (one family), incorrect.
Use groupings to decide whether Lesson 3.11 opens with a brief recap of the
general-solution structure.
\end{skillbox}

\vspace{0.08in}

%----------------------------------------------------------------------
% Reinforcement & Extension
%----------------------------------------------------------------------
\begin{skillbox}[Reinforcement \& Extension]{goldbox}
\textbf{Homework (Lesson 3.10):} 5--6 problems: solving basic trig equations
for all solutions in a given interval, writing the general solution, one
contextual domain problem with interpretation, 2 AP-style MC items, and an
extension asking students to solve $2\sin^2\theta - \sin\theta - 1 = 0$ via
quadratic substitution $u = \sin\theta$.

\textbf{Preview of Lesson 3.11 --- Secant, Cosecant, and Cotangent:}
Ask students to consider: \textit{``We have three more trig functions ---
$\sec$, $\csc$, and $\cot$ --- defined as the reciprocals of cosine, sine,
and tangent. How would you expect to solve an equation like $\sec\theta = 2$?
What period would you predict?''} Lesson 3.11 extends today's skills to these
reciprocal functions.
\end{skillbox}

\end{document}
